\documentclass[border=2pt]{standalone}
\usepackage{amsmath, amsthm, amsfonts}
\usepackage{tikz-cd, kotex}
\usepackage{quiver}
\usepackage{Operators}
\usepackage{palette}
\usepackage{diagram-style}
\usepackage{blowup-surface}
\begin{document}

%%FIG the curves over the lines x=epsilon tend to E and the fiber over the pole (Research_Note_2026_11_19-1.svg)
\begin{tikzpicture}[line cap=round, line join=round]
  %% left: the chart (s,t) of X around the point (q,l_0); E={t=0}, F_{l_0}={s=0}
  \foreach \a in {-1.8,-.9,.9,1.8}{\draw[black!25, line width=.5pt] (\a,-2.1)--(\a,2.1);}
  \foreach \e/\w/\c in {1/1.1pt/black, .4/.7pt/black!50, .12/.7pt/black!50}{
    \draw[\c, line width=\w] plot[domain={\e/2.1}:2.1, samples=70] (\x,{\e/\x});
    \draw[\c, line width=\w] plot[domain=-2.1:{-\e/2.1}, samples=70] (\x,{\e/\x});
  }
  \draw[accent1, line width=1.2pt] (-2.3,0)--(2.3,0);
  \node[accent1, right] at (2.3,0) {$E$};
  \draw[accent2, line width=1.2pt] (0,-2.3)--(0,2.3);
  \node[accent2, above] at (0,2.3) {$F_{\ell_0}$};
  \fill (0,0) circle[radius=1.6pt];
  \node[below right, font=\scriptsize] at (0,0) {$(q,\ell_0)$};
  \node[right, font=\scriptsize] at (1.05,1.2) {$\mathrm{s}\mathrm{t}=\epsilon$};
  \node at (0,-2.9) {$X$};
  %% the blowup map
  \draw[thick, -{stealth}] (3.4,0)--(4.6,0);
  \node[above] at (4,0) {$\pi$};
  %% right: the plane around q; lines through q, l_0={x=0}, L_epsilon={x=epsilon}
  \begin{scope}[shift={(7.6,0)}]
    \draw[black!25, line width=.5pt] (-1.89,-2.1)--(1.89,2.1);
    \draw[black!25, line width=.5pt] (1.89,-2.1)--(-1.89,2.1);
    \draw[black!25, line width=.5pt] (-2.1,-1.167)--(2.1,1.167);
    \draw[black!25, line width=.5pt] (2.1,-1.167)--(-2.1,1.167);
    \foreach \e/\w/\c in {1/1.1pt/black, .4/.7pt/black!50, .12/.7pt/black!50}{
      \draw[\c, line width=\w] (\e,-2.1)--(\e,2.1);
    }
    \draw[accent2, line width=1.2pt] (0,-2.3)--(0,2.3);
    \node[accent2, above] at (0,2.3) {$\ell_0$};
    \fill[accent1] (0,0) circle[radius=2pt];
    \node[accent1, below left] at (0,0) {$q$};
    \node[right, font=\scriptsize] at (1,1.75) {$L_\epsilon$};
    \node at (0,-2.9) {$\mathbb{P}^2$};
  \end{scope}
\end{tikzpicture}

%%FIG the two kinds of stable maps sending the marking x_n into E (Research_Note_2026_11_19-2.svg)
\begin{tikzpicture}[line cap=round, line join=round]
  %% left: the component carrying x_n is not contained in E and meets E at f(x_n)
  \draw[accent1, line width=1.2pt] (-2,0)--(2,0);
  \node[accent1, right] at (2,0) {$E$};
  \draw[line width=1.5pt] (-1.3,-1.4) .. controls (-1,-.6) and (-.4,-.3) .. (0,0)
    .. controls (.4,.3) and (1,.5) .. (1.4,1.4);
  \fill (0,0) circle[radius=2.4pt];
  \node[below right] at (0,0) {$f(x_n)$};
  %% right: the component carrying x_n is mapped into E; the other components cross E
  \begin{scope}[shift={(6.2,0)}]
    \draw[black!50, line width=1.5pt] (-1.6,-1.4) .. controls (-1.4,-.6) and (-1.3,-.3) .. (-1.2,0)
      .. controls (-1.1,.3) and (-.9,.8) .. (-.4,1.4);
    \draw[black!50, line width=1.5pt] (1.7,-1.4) .. controls (1.5,-.6) and (1.4,-.3) .. (1.3,0)
      .. controls (1.2,.3) and (1.3,.9) .. (1.8,1.4);
    \draw[accent1, line width=1.2pt] (-2.5,0)--(2.5,0);
    \node[accent1, right] at (2.5,0) {$E$};
    \draw[line width=1.5pt] (-1.7,0)--(1.7,0);
    \fill (.1,0) circle[radius=2.4pt];
    \node[below] at (.1,0) {$f(x_n)$};
  \end{scope}
\end{tikzpicture}

%%FIG the blowup X of the plane at q, on which the lines through q are pulled apart along E (Research_Note_2026_11_19-3.svg)
\begin{tikzpicture}[line cap=round, line join=round]
  \BlowupSurface
  \node[accent1, left] at ($(\bpt{-\Bphimax}{0})+(-.5,0)$) {$E$};
  \node at (0,-2.45) {$X$};
  \draw[thick, ->] (4.75,0)--(5.55,0);
  \node[above] at (5.15,0) {$\pi$};
  \begin{scope}[shift={(6.75,0)}]
    \BlowupDisc
    \fill[accent1] (0,0) circle[radius=2pt];
    \node[accent1, left] at (-.08,-.12) {$q$};
    \node at (0,-2.45) {$\mathbb{P}^2$};
  \end{scope}
\end{tikzpicture}

%%FIG the ruling rho, whose fiber over l is the strict transform of the line l (Research_Note_2026_11_19-4.svg)
\begin{tikzpicture}[line cap=round, line join=round]
  \BlowupMark{-2}{accent2, line width=1.3pt}
  \BlowupSurface
  \node[accent1, left] at ($(\bpt{-\Bphimax}{0})+(-.5,0)$) {$E$};
  \fill[accent2] (\bpt{bang(-2)}{0}) circle[radius=2pt];
  \node[accent2, above] at (\bpt{bang(-2)}{\BR}) {$\rho^{-1}(\ell)$};
  %% the base of rho
  \draw[black!50, line width=.45pt, densely dashed] (\bpt{bang(-2)}{0})--({bu(bang(-2))},-2.95);
  \draw[line width=1.1pt] ($(\bpt{-\Bphimax}{0})+(-.5,-2.95)$)--($(\bpt{\Bphimax}{0})+(.5,-2.95)$);
  \node[left] at ($(\bpt{-\Bphimax}{0})+(-.5,-2.95)$) {$\mathbb{P}^1_q$};
  \fill[accent2] ({bu(bang(-2))},-2.95) circle[radius=2pt];
  \node[accent2, below] at ({bu(bang(-2))},-2.95) {$\ell$};
  \draw[thick, ->] (1.2,-2.0)--(1.2,-2.75);
  \node[right] at (1.2,-2.35) {$\rho$};
  \draw[thick, ->] (4.75,0)--(5.55,0);
  \node[above] at (5.15,0) {$\pi$};
  \begin{scope}[shift={(6.75,0)}]
    \BlowupDisc
    \fill[accent1] (0,0) circle[radius=2pt];
    \node[accent1, left] at (-.08,-.12) {$q$};
    \node[accent2, above] at (\dpt{bang(-2)}{\BR}) {$\ell$};
    \node at (0,-2.05) {$\mathbb{P}^2$};
  \end{scope}
\end{tikzpicture}

%%FIG the blowup, the ruling and the projection from q (Research_Note_2026_11_19-5.svg)
\begin{tikzcd}
E \arrow[r] \arrow[d, hook] \arrow[dd, bend right=60, "\cong"'] & \{q\} \arrow[d, hook] \\
X \arrow[r, "\pi"] \arrow[d, "\rho"] & \mathbb{P}^2 \arrow[ld, dashed] \\
\mathbb{P}^1_q &
\end{tikzcd}

%%FIG curves of the classes H, F and E on X (Research_Note_2026_11_19-6.svg)
\begin{tikzpicture}[line cap=round, line join=round]
  \BlowupMark{-2}{accent2, line width=1.3pt}
  %% the line L: at distance \Ld from q, with normal direction \Lnu
  \def\Ld{1.0}\def\Lnu{-20}
  \pgfmathsetmacro{\Llo}{\Lnu-acos(\Ld/\BR)}\pgfmathsetmacro{\Lhi}{\Lnu+acos(\Ld/\BR)}
  \renewcommand{\BlowupSliceHook}[2]{%
    \pgfmathsetmacro{\La}{max(#1,\Llo)}\pgfmathsetmacro{\Lb}{min(#2,\Lhi)}%
    \pgfmathparse{\La<\Lb}\ifnum\pgfmathresult=1
      \draw[line width=1.3pt] (\bpt{\La}{\Ld/cos(\La-\Lnu)})--(\bpt{\Lb}{\Ld/cos(\Lb-\Lnu)});
    \fi}
  \BlowupSurface
  \node[accent1, left] at ($(\bpt{-\Bphimax}{0})+(-.5,0)$) {$E$};
  \fill[accent2] (\bpt{bang(-2)}{0}) circle[radius=2pt];
  \fill (\bpt{bang(-2)}{\Ld/cos(bang(-2)-\Lnu)}) circle[radius=2pt];
  \node[accent2, above] at (\bpt{bang(-2)}{\BR}) {$F$};
  \node[above] at (\bpt{\Lhi}{\BR}) {$H$};
  \draw[thick, ->] (4.75,0)--(5.55,0);
  \node[above] at (5.15,0) {$\pi$};
  \begin{scope}[shift={(6.75,0)}]
    \BlowupDisc
    \draw[line width=1.3pt] (\dpt{\Llo}{\BR})--(\dpt{\Lhi}{\BR});
    \fill (\dpt{bang(-2)}{\Ld/cos(bang(-2)-\Lnu)}) circle[radius=2pt];
    \fill[accent1] (0,0) circle[radius=2pt];
    \node[accent1, left] at (-.08,-.12) {$q$};
    \node[accent2, above] at (\dpt{bang(-2)}{\BR}) {$\ell$};
    \node[above right] at (\dpt{\Lhi}{\BR}) {$L$};
    \node at (0,-2.05) {$\mathbb{P}^2$};
  \end{scope}
\end{tikzpicture}

%%FIG the tangent directions of E and of the fiber at a point of E (Research_Note_2026_11_19-7.svg)
\begin{tikzpicture}[line cap=round, line join=round]
  \BlowupMark{-2}{accent2, line width=.8pt}
  \BlowupSurface
  \node[accent1, left] at ($(\bpt{-\Bphimax}{0})+(-.5,0)$) {$E$};
  \draw[accent1, line width=1.6pt, ->] (\bpt{bang(-2)}{0})--++(.95,0);
  \draw[accent2, line width=1.6pt, ->] (\bpt{bang(-2)}{0})--(\bpt{bang(-2)}{.95});
  \fill (\bpt{bang(-2)}{0}) circle[radius=2pt];
  \node[accent1, below, fill=white, inner sep=1.5pt] at ($(\bpt{bang(-2)}{0})+(.72,-.12)$) {$T_E$};
  \node[accent2, left, fill=white, inner sep=1.5pt] at ($(\bpt{bang(-2)}{.8})+(-.12,0)$) {$T_{X/\mathbb{P}^1}$};
\end{tikzpicture}

%%FIG the relative tangent sequence and the normal sequence of E, restricted to E (Research_Note_2026_11_19-8.svg)
\begin{tikzcd}
 & & 0 \arrow[d] & & \\
 & & T_E \arrow[d] \arrow[rd, "\cong"] & & \\
0 \arrow[r] & T_{X/\mathbb{P}^1}\vert_E \arrow[r] \arrow[rd, "\cong"'] & T_X\vert_E \arrow[r, "d\rho"] \arrow[d] & \rho^\ast T_{\mathbb{P}^1}\vert_E \arrow[r] & 0 \\
 & & N_{E/X} \arrow[d] & & \\
 & & 0 & &
\end{tikzcd}

\end{document}
